% ============================================================
\chapter{Chaos --- Definitions and Examples}
\label{ch:chaos}
% ============================================================

\section{Introduction}

In 1963, Edward Lorenz, a meteorologist at MIT, discovered by accident that his simplified model of atmospheric convection produced radically different results depending on whether an initial condition was rounded to three or six decimal places. This phenomenon --- \emph{sensitivity to initial conditions} --- would become the symbol of deterministic chaos, popularized under the name ``butterfly effect.'' But what exactly is chaos? The question is subtler than it appears. Robert Devaney proposed a rigorous definition in 1989 combining three ingredients: sensitivity to initial conditions, topological transitivity, and density of periodic orbits. Li and Yorke, in their celebrated 1975 paper ``Period Three Implies Chaos,'' had given yet another. This chapter explores these definitions, their relationships, and the fundamental properties of chaotic systems --- those deterministic systems that, paradoxically, produce the appearance of randomness.

\section{Sensitivity to Initial Conditions}

\begin{definition}[Sensitivity to initial conditions]
\index{sensitivity to initial conditions}
A dynamical system $f : X \to X$ on a metric space $(X, d)$ has
\textbf{sensitivity to initial conditions} (SIC) if there exists $\delta > 0$
such that for every $x \in X$ and every $\epsilon > 0$, there exists $y$ with
$d(y, x) < \epsilon$ and $n > 0$ such that
\[
    d(f^n(x), f^n(y)) > \delta.
\]
\end{definition}

\begin{intuition}
Sensitivity to initial conditions is often called the \emph{butterfly effect}:
an infinitesimal perturbation can lead to radically different trajectories.
This is the fundamental discovery of Lorenz (1963).
\end{intuition}

\begin{remark}
Sensitivity alone does not define chaos. The map $x \mapsto 2x$ on $\R$ is
sensitive to initial conditions but is not chaotic --- orbits simply escape
to infinity without any recurrence.
\end{remark}

\section{Devaney Chaos}

\begin{definition}[Devaney chaos]
\label{def:devaney}
\index{chaos!Devaney}
A continuous map $f : X \to X$ on a metric space $(X, d)$ is
\textbf{chaotic in the sense of Devaney} if:
\begin{enumerate}
    \item \textbf{Topological transitivity}: for all non-empty open sets
    $U, V \subseteq X$, there exists $n \in \N$ with $f^n(U) \cap V \neq \emptyset$.
    \item \textbf{Dense periodic orbits}: the set of periodic points is dense in $X$.
    \item \textbf{Sensitivity to initial conditions}.
\end{enumerate}
\end{definition}

\begin{theorem}[Banks et al., 1992]
\index{Banks theorem}
If $f : X \to X$ is topologically transitive with dense periodic orbits on an
infinite metric space, then $f$ automatically has sensitivity to initial
conditions. Thus condition (3) is \textbf{redundant} in Devaney's definition.
\end{theorem}

\begin{proof}[Proof sketch]
Since periodic orbits are dense and $X$ is infinite, there exist periodic
points $p, q$ with $d(p, q) > 0$. Setting $\delta = d(p, q)/4$, for any
$x \in X$ and $\epsilon > 0$, transitivity yields iterates that visit
neighbourhoods of both $p$ and $q$, producing the required separation.
\end{proof}

\section{Li-Yorke Chaos}

\begin{definition}[Li-Yorke pair]
\index{chaos!Li-Yorke}
Two points $x, y$ form a \textbf{Li-Yorke pair} if
\[
    \limsup_{n \to \infty} d(f^n(x), f^n(y)) > 0 \quad \text{and} \quad
    \liminf_{n \to \infty} d(f^n(x), f^n(y)) = 0.
\]
\end{definition}

\begin{definition}[Li-Yorke chaos]
The system is \textbf{Li-Yorke chaotic} if there exists an uncountable
\textbf{scrambled set} $S$ such that every pair of distinct points in $S$ is
a Li-Yorke pair.
\end{definition}

\begin{theorem}[Li-Yorke, 1975]
\label{thm:li-yorke}
Let $f : [0,1] \to [0,1]$ be continuous. If $f$ has a periodic orbit of
period~3, then $f$ has periodic orbits of every period $n \in \N^*$, and $f$
is Li-Yorke chaotic.
\end{theorem}

\begin{remark}
The famous statement ``Period three implies chaos'' summarises this theorem.
More generally, Sharkovskii's ordering determines which periods force which
others.
\end{remark}

\section{Sharkovskii's Theorem}

\begin{theorem}[Sharkovskii, 1964]
\index{Sharkovskii's theorem}
Define the Sharkovskii ordering on $\N^*$:
\[
    3 \triangleright 5 \triangleright 7 \triangleright \cdots \triangleright
    2\cdot 3 \triangleright 2\cdot 5 \triangleright \cdots \triangleright
    4\cdot 3 \triangleright 4\cdot 5 \triangleright \cdots \triangleright
    \cdots \triangleright 8 \triangleright 4 \triangleright 2 \triangleright 1.
\]
If $f : \R \to \R$ is continuous and has a periodic orbit of period $m$, then
$f$ has periodic orbits of every period $n$ with $m \triangleright n$.
Moreover, for each $m$ there exists a continuous map with a period-$m$ orbit
but no period-$n$ orbit for any $n \triangleright m$.
\end{theorem}

\section{Fundamental Examples}

\subsection{The shift map}

\begin{definition}[Sequence space and shift]
\index{shift map}
The space of binary sequences is
$\Sigma_2 = \{0,1\}^\N$
with metric $d(s,t) = \sum_{i=0}^{\infty} \abs{s_i - t_i}/2^i$. The
\textbf{shift} $\sigma : \Sigma_2 \to \Sigma_2$ is defined by
$\sigma(s_0, s_1, s_2, \ldots) = (s_1, s_2, s_3, \ldots)$.
\end{definition}

\begin{theorem}
The shift $\sigma$ on $\Sigma_2$ is Devaney chaotic:
\begin{enumerate}
    \item It is topologically transitive.
    \item Periodic points are dense.
    \item It has sensitivity to initial conditions with $\delta = 1$.
\end{enumerate}
\end{theorem}

\subsection{The tent map}

\begin{definition}[Tent map]
\index{tent map}
The \textbf{tent map} is $T : [0,1] \to [0,1]$ defined by
\[
    T(x) = \begin{cases} 2x & \text{if } 0 \leq x \leq 1/2, \\
    2(1-x) & \text{if } 1/2 < x \leq 1. \end{cases}
\]
\end{definition}

\begin{proposition}
The tent map is semi-conjugate to the shift on $\Sigma_2$ and is Devaney
chaotic.
\end{proposition}

\subsection{The logistic map at $r=4$}

\begin{theorem}
The logistic map $f_4(x) = 4x(1-x)$ on $[0,1]$ is topologically conjugate
to the tent map via $h(x) = \frac{2}{\pi}\arcsin(\sqrt{x})$. It is therefore
Devaney chaotic.
\end{theorem}

\section{Smale's Horseshoe}

\begin{definition}[Smale horseshoe]
\index{Smale horseshoe}
The \textbf{Smale horseshoe} is a diffeomorphism of the plane that stretches
a square, folds it into a horseshoe shape, and maps it back into the original
square. The maximal invariant set is a Cantor set on which the dynamics is
conjugate to the shift on $\Sigma_2$.
\end{definition}

\begin{theorem}[Smale, 1967]
On the invariant set of the horseshoe, the dynamics is:
\begin{enumerate}
    \item Devaney chaotic;
    \item uniformly hyperbolic;
    \item structurally stable (persists under small perturbations).
\end{enumerate}
\end{theorem}

\section{Routes to Chaos}

\begin{definition}[Period-doubling cascade]
\index{period-doubling cascade}
The \textbf{Feigenbaum cascade} is an infinite sequence of period-doubling
bifurcations $1 \to 2 \to 4 \to 8 \to \cdots$ leading to chaos.
\end{definition}

\begin{theorem}[Feigenbaum universality]
\index{Feigenbaum universality}
The bifurcation values $\mu_n$ satisfy
\[
    \lim_{n \to \infty} \frac{\mu_n - \mu_{n-1}}{\mu_{n+1} - \mu_n}
    = \delta_F \approx 4.6692\ldots
\]
The constant $\delta_F$ is \textbf{universal}: it does not depend on the
particular family of maps.
\end{theorem}

\begin{keyformulas}[title=\textbf{Feigenbaum Constants}]
\begin{align*}
    \delta_F &\approx 4.669\,201\,609\ldots
    \quad \text{(parameter space scaling ratio)} \\
    \alpha_F &\approx -2.502\,907\,875\ldots
    \quad \text{(spatial scaling factor)}
\end{align*}
\end{keyformulas}

Other routes to chaos:
\begin{itemize}
    \item \textbf{Intermittency} (Pomeau-Manneville): alternation of laminar
    phases and chaotic bursts.
    \item \textbf{Quasi-periodicity} (Ruelle-Takens): destruction of an
    invariant torus.
    \item \textbf{Crisis}: collision of a chaotic attractor with a repeller.
\end{itemize}

\section{Chaos in Continuous Systems}

\begin{attention}[title=\textbf{Minimum dimension for continuous chaos}]
An autonomous continuous system $\dot{x} = f(x)$ can only be chaotic in
dimension $n \geq 3$. In dimension 2, the Poincar\'e-Bendixson theorem
forbids chaos.
\end{attention}

\begin{example}[Mel'nikov's theorem]
\index{Mel'nikov theorem}
Mel'nikov's theorem detects chaos in perturbed Hamiltonian systems. For the
forced pendulum $\ddot{x} + \sin x = \epsilon(\gamma\cos\omega t - \delta\dot{x})$,
the Mel'nikov function predicts the existence of transverse homoclinic orbits
(hence a Smale horseshoe) for sufficiently small $\epsilon$.
\end{example}

\begin{codeblock}[title=\textbf{Demonstrating sensitivity to initial conditions}]
\begin{minted}{python}
import numpy as np
from scipy.integrate import solve_ivp
import matplotlib.pyplot as plt

def lorenz(t, state, sigma=10, rho=28, beta=8/3):
    x, y, z = state
    return [sigma*(y - x), rho*x - y - x*z, x*y - beta*z]

x0_1 = [1.0, 1.0, 1.0]
x0_2 = [1.0 + 1e-10, 1.0, 1.0]
t_span, t_eval = [0, 50], np.linspace(0, 50, 5000)
sol1 = solve_ivp(lorenz, t_span, x0_1, t_eval=t_eval, max_step=0.01)
sol2 = solve_ivp(lorenz, t_span, x0_2, t_eval=t_eval, max_step=0.01)

fig, (ax1, ax2) = plt.subplots(2, 1, figsize=(12, 8))
ax1.plot(sol1.t, sol1.y[0], 'b-', lw=0.5, label='IC 1')
ax1.plot(sol2.t, sol2.y[0], 'r-', lw=0.5, label='IC 2')
ax1.legend(); ax1.set_title('Sensitivity to initial conditions')

diff = np.sqrt(np.sum((sol1.y - sol2.y)**2, axis=0))
ax2.semilogy(sol1.t, diff, 'k-', lw=0.5)
ax2.set_ylabel(r'$\|\Delta x(t)\|$')
ax2.set_title('Exponential divergence')
plt.tight_layout(); plt.savefig("sensitivity.pdf")
\end{minted}
\end{codeblock}

\section{Exercises}

\begin{exercise}[Shift and chaos]
Show that the shift $\sigma$ on $\Sigma_2$ is topologically transitive by
explicitly constructing a dense orbit.
\end{exercise}

\begin{exercise}[Period three]
Show that $f(x) = x^2 - 7/4$ on $\R$ has a period-3 orbit and deduce that it
is Li-Yorke chaotic.
\end{exercise}

\begin{exercise}[Tent map periodic points]
Find all periodic points of period $\leq 3$ of the tent map and verify that
they are all unstable.
\end{exercise}

\begin{exercise}[Logistic map]
For $f_r(x) = rx(1-x)$, numerically demonstrate the period-doubling cascade
and estimate the Feigenbaum constant.
\end{exercise}

\begin{exercise}[Horseshoe]
Describe geometrically the construction of the Smale horseshoe and show that
the invariant set is a Cantor product.
\end{exercise}

\begin{exercise}[Three-dimensional chaos]
Why can the R\"ossler system $\dot{x} = -y - z$, $\dot{y} = x + ay$,
$\dot{z} = b + z(x - c)$ be chaotic despite having only three equations?
Verify numerically for $a = 0.2$, $b = 0.2$, $c = 5.7$.
\end{exercise}
